\chapter{Infrastruttura comune}

L'infrastruttura di supporto \`e condivisa da tutti i modelli e comprende il dataset PyTorch, il pipeline di training, il sistema di esperimenti, le loss function e le utilit\`a di valutazione e inferenza.

\section{Dataset e Data Loading}

La classe \texttt{ImageEnhancementDataset} estende \texttt{torch.utils.data.Dataset} e gestisce il caricamento delle coppie (immagine degradata, immagine pulita) con estrazione di patch casuali.

\subsection{Estrazione delle patch}

Ogni immagine viene suddivisa in patch quadrate di dimensione configurabile (tipicamente $128\!\times\!128$). Per ogni immagine vengono estratte \texttt{patches\_per\_image} patch in posizioni casuali, cos\`i che la lunghezza effettiva del dataset \`e:
$$
|\mathcal{D}| = N_{\text{immagini}} \times \text{patches\_per\_image}
$$

Le immagini vengono mantenute in cache tramite \texttt{@lru\_cache(maxsize=64)} per ridurre gli accessi a disco.

\subsection{Data augmentation}

Durante il training vengono applicate trasformazioni geometriche tramite la libreria Albumentations:
\begin{itemize}
    \item \texttt{HorizontalFlip} (probabilit\`a 0.5)
    \item \texttt{VerticalFlip} (probabilit\`a 0.5)
    \item \texttt{RandomRotate90} (probabilit\`a 0.5)
\end{itemize}

Le augmentation vengono applicate congiuntamente alla coppia degradata/pulita tramite il meccanismo \texttt{additional\_targets} di Albumentations. La normalizzazione mappa i pixel da $[0, 255]$ a $[-1, 1]$ tramite $x' = (x/255 - 0.5) / 0.5$. In validazione si applica solo la normalizzazione.

\subsection{DataLoader}

La funzione \texttt{get\_dataloaders} crea i dataloader di training e validazione con le seguenti impostazioni:
\begin{itemize}
    \item \texttt{pin\_memory=True} per trasferimenti GPU asincroni
    \item \texttt{persistent\_workers=True} per evitare il respawn dei worker
    \item \texttt{drop\_last=True} nel training per batch di dimensione costante
    \item Worker init function che disabilita il threading OpenCV e imposta seed NumPy per-worker
\end{itemize}

\section{Funzioni di Loss}

Sono state implementate diverse loss function, dalle pi\`u semplici a combinazioni pi\`u elaborate.

\subsection{Combined Loss (L1 + SSIM)}

La loss combinata \`e definita come:
$$
\mathcal{L}_{\text{combined}} = \alpha \cdot \mathcal{L}_{\text{L1}} + \beta \cdot (1 - \text{SSIM})
$$

dove $\alpha = 0.84$ e $\beta = 0.16$ sono i pesi di default. Per il calcolo dell'SSIM le immagini vengono convertite da $[-1,1]$ a $[0,1]$. Per stabilizzare il calcolo in mixed precision (FP16), vengono applicati un epsilon di clamping e una piccola iniezione di rumore gaussiano.

\subsection{Perceptual Loss (VGG)}

La \texttt{VGGPerceptualLoss} utilizza una rete VGG-16 pre-addestrata su ImageNet come estrattore di feature. Le feature vengono estratte a livelli configurabili (\texttt{relu2\_2}, \texttt{relu3\_3} di default) e la loss \`e calcolata come L1 pesata tra le feature map:
$$
\mathcal{L}_{\text{perceptual}} = \sum_{l \in \text{layers}} w_l \cdot \|\phi_l(\hat{x}) - \phi_l(x)\|_1
$$

dove $\phi_l$ indica le feature al livello $l$ della VGG-16. I parametri della VGG sono congelati e gli input vengono normalizzati con le statistiche di ImageNet.

\subsection{Combined Perceptual Loss}

La \texttt{CombinedPerceptualLoss} unifica le tre componenti:
$$
\mathcal{L} = \alpha \cdot \mathcal{L}_{\text{L1}} + \beta \cdot (1 - \text{SSIM}) + \gamma \cdot \mathcal{L}_{\text{perceptual}}
$$

Con i pesi tipici $\alpha = 0.6$, $\beta = 0.25$, $\gamma = 0.15$. Quando $\gamma = 0$ la rete VGG non viene caricata in memoria.

\section{Sistema di Esperimenti}

Il modulo \texttt{experiment.py} gestisce la creazione, il salvataggio e la ripresa degli esperimenti.

\subsection{Struttura delle directory}

Ogni esperimento viene salvato in una directory con timestamp:
\begin{verbatim}
experiments/{modello}/{degradazione}/{timestamp}_{nome}/
    checkpoints/    # Pesi del modello
    samples/        # Campioni visivi
    logs/           # Log TensorBoard
    config.json     # Configurazione completa
    history.json    # Metriche per epoch
\end{verbatim}

\subsection{Ripresa e fork}

Il sistema supporta la ripresa automatica dell'ultimo esperimento. Se la configurazione \`e cambiata rispetto all'esperimento precedente, viene eseguito un \textit{fork}: si crea un nuovo esperimento con i propri metadati, copiando il miglior checkpoint e la history dal padre. Questo consente di esplorare variazioni di iperparametri mantenendo la tracciabilit\`a.

La funzione \texttt{compare\_configs} confronta le configurazioni ignorando campi non significativi come \texttt{resume\_from\_checkpoint}, \texttt{num\_epochs} e \texttt{device}.

\section{Checkpointing}

I checkpoint salvano lo stato completo dell'addestramento:
\begin{itemize}
    \item Pesi del modello (\texttt{model\_state\_dict})
    \item Stato dell'ottimizzatore e dello scheduler
    \item Epoch corrente e metriche
    \item Metadati opzionali (configurazione, ecc.)
\end{itemize}

Al caricamento, viene eseguita una validazione di integrit\`a: tutti i tensori floating-point vengono controllati per la presenza di valori \texttt{NaN} o \texttt{Inf}, con errore esplicito che indica le chiavi corrotte.

In caso di errore OOM durante il training, viene salvato un checkpoint di emergenza. Se il file \texttt{best\_model.pth} manca, il sistema lo ricrea a partire dal checkpoint di emergenza o periodico pi\`u recente.

\section{Notifiche Telegram}

Il modulo \texttt{telegram\_notifier.py} invia notifiche sullo stato del training tramite l'API Telegram Bot. Vengono notificati: inizio training, aggiornamenti periodici sulle epoche, completamento, early stopping e errori OOM. Le credenziali vengono lette da un file \texttt{.env} (\texttt{TELEGRAM\_BOT\_TOKEN} e \texttt{TELEGRAM\_CHAT\_ID}).
